\chapter{Theory}\label{chap: Theory}


\section{Description of turbulent flows}\label{sec: Description of turbulent flows}

In turbulent flows, structures in a multitude of scales emerge.
The minimum scale of the elements is limited by the viscosity of the fluid, since the smallest scales are dissipated by viscosity and disappear.
Thus, the viscous units, or plus units are introduced. 

\subsection{Wall shear stress}\label{subsec: wall shear stress}

The wall shear stress, which is responsible for the drag, is defined as 
\begin{equation}
    \tau_w = \mu\frac{dU}{dy}
\end{equation}
and is solely based on viscosity (\cite{pope_turbulent_2015}).
The viscosity is dominant at the wall, hence the wall shear stress is the only stress active at the wall and it keeps being dominant in the near wall region.
With a growing distance to the wall the influence of viscosity decreases, while the importance of the Reynolds stresses, described in section \ref{subsec: Reynolds stress} increases.


\subsection{Friction velocity}\label{subsec: Friction velocity}

Using the wall shear stress and the density $\rho$ the friction velocity is defined as
\begin{equation}
    u_\tau = \sqrt{\frac{\tau_w}{\rho}}
\end{equation}
The friction velocity $u_\tau$ is a one of the viscous scales, which are used to a
\subsection{Reynolds stresses}\label{subsec: Reynolds stress}


%\section{Numerics}\label{sec: Numerics}

\section{Roughness}\label{sec: Roughness}

\subsection{Roughness parameters}\label{subsec: Roughness parameters}

To be able to describe rough surfaces, a multitude of roughness parameters exist, as described by \cite{chung_predicting_2021}.
The considered roughness parameters in this study are:

\paragraph{Average roughness height}

The average roughness height is the mean deviation  of the surface from the virtual origin $y=0$
\begin{equation}\label{eq: avg roughness height}
    S_a = \sum_{i,j}^{M,N} |h_{i,j}| / (MN)
\end{equation}

\paragraph{Root-mean-squeare roughness height}

\begin{equation}
    S_q = \sqrt{\sum_{i,j}^{M,N} h_{i,j}^{2}} / (MN)
\end{equation}

\paragraph{Skewness}

\begin{equation}
    Sk = \sum_{i,j}^{M,N} h^{3}_{i,j} /(MNS_{q}^{3})
\end{equation}

\paragraph{Maximum peak-to-valley height}
The maximum peak-to-valley height is the difference between the highest and lowest point of the rough surface
\begin{equation}
    S_{Z_{max}} = max(h_{i,j}) - min(h_{i,j})
\end{equation}

\paragraph{Effective slope}
The effective slope is the absolute gradient of the rough surface as described by \cite{napoli_effect_2008}.
\begin{equation}
    ES = \frac{1}{L_x L_y} \int_0^{L_x} \int_0^{L_y} \left| \frac{\mathrm{d}h}{\mathrm{d}x} \right|_{i,j} \mathrm{d}x \, \mathrm{d}y
\end{equation}
\paragraph{Mean peak-to-valley height}
The calculation of the mean peak to valley height is done according to \cite{thakkar_investigation_2017}.
The method is dividing the surface into tiles and calculating the difference between the highest and lowest point for each tile.
The mean of the differences results in  the mean peak-to-valley height $k$, which is also used as the measurement for roughness height in this thesis, according to \cite{de_maio_direct_2023}.
For this study the rough surface is divided into 25 rectangular tiles of equal size.
Equal size means in this context, that the same number of grid points are used for the length and width of the tiles.


\subsection{Virtual origin}\label{subsec: Virtual origin}

The virtual origin is a reference plane, which is taken as the basis for measuring the wall distance $y$ (\cite{chung_predicting_2021}).
In a smooth wall case, there is no virtual origin, since all data can be calculated from the wall itself, hence the point where the no-slip condition is active (\cite{chan_systematic_2015}), which leads to consistent wall distances for every point in the field.\\
Yet, in a rough wall case, this is not as straightforward.
Due to wall roughness, the distance from e.g. the pipe center to the wall varies continuously over the whole flow field.
To be able to compute consistent data, a virtual origin has to be defined. 
There are many different approaches to defining the virtual origin.
Notable methods are the usage of the mean momentum absorption plane by \cite{thom_momentum_1971} or the determination of the position of the mean streamwise velocity being zero, since it is assumed, that this is the equivalent position of the no slip condition in a smooth wall case.
A third approach is to take the mean radius of the rough surface for the virtual origin.\\
In the investigation by \cite{chan_systematic_2015}, all three of the approaches are compared. 
The approach of using the mean radius for the virtual origin seems to offer the best correlation between the considered smooth and rough cases in this study.
For that reason, this is the used method in this study.


